%% guide-bac-virage-releve — coupe transversale d'un virage relevé de circuit de karting :
%% sol horizontal, piste inclinée du dévers de 32° (donnée, cotée par un arc), kart
%% schématisé posé perpendiculairement à la piste, verticale passant par le point de
%% contact I.
%% \rappel cache le cône de frottement en I : la normale à la piste, les deux génératrices
%% à phi = 36,9° de part et d'autre, l'aplat clair du cône et la conclusion (la résultante
%% du contact reste dans le cône, donc pas de glissement). Le cône est dessiné AVANT le
%% kart pour passer derrière lui, la légende après.
\documentclass[tikz,border=4pt]{standalone}
\usepackage{liaisons}
\begin{document}
\begin{tikzpicture}[x=1cm,y=1cm,
  piste/.style={line width=0.9pt, line join=round, draw=black!70, fill=black!10},
  axe/.style={line width=0.5pt, draw=black!60,
              dash pattern=on 8pt off 2.5pt on 1.5pt off 2.5pt},
  cone/.style={line width=0.9pt, draw=solideA},
  etiq/.style={font=\small, inner sep=2pt}]

\def\dev{32}                       % dévers de la piste
\def\phif{36.9}                    % angle de frottement
\coordinate (T) at (2.40,0.50);    % pied du dévers
\coordinate (P) at (6.98,3.36);    % haut de la piste : T + 5,40 cm à 32°
\coordinate (I) at (4.859,2.037);  % point de contact, à 2,90 cm de T le long de la piste

%% ---------------- le sol horizontal et le talus de la piste -----------------------------
\hachures[draw=black!70]{(0.30,0.18) rectangle (7.40,0.50)}
\draw[piste] (T) -- (P) -- (6.98,0.50) -- cycle;
\begin{scope}[shift={(T)}, rotate=\dev]
  \draw[line width=0.9pt, draw=black!75, fill=black!25] (0,-0.16) rectangle (5.40,0);
\end{scope}
\node[etiq, anchor=west, text=black!75, align=left] at (7.10,3.30) {piste relevée};
\node[etiq, anchor=north west, text=black!75] at (0.30,0.14) {sol horizontal};

%% ---------------- le dévers, coté (donnée de l'énoncé) ----------------------------------
\draw[line width=0.7pt, draw=solideD] ([shift=(0:1.55)]T) arc (0:\dev:1.55);
\node[etiq, anchor=west, text=solideD] at (4.05,0.94) {dévers $32^\circ$};

%% ---------------- la réponse (1) : le cône de frottement, derrière le kart ---------------
\rappel{%
  \fill[solideA!10] (I) -- ([shift=({90+\dev-\phif}:1.95)]I)
     arc ({90+\dev-\phif}:{90+\dev+\phif}:1.95) -- cycle;
  \draw[cone] (I) -- ([shift=({90+\dev-\phif}:1.95)]I);
  \draw[cone] (I) -- ([shift=({90+\dev+\phif}:1.95)]I);
  \draw[axe, draw=solideA!80] (I) -- ([shift=({90+\dev}:2.15)]I);}

%% ---------------- la verticale du point de contact (par-dessus l'aplat du cône) ----------
\draw[axe] (4.859,0.28) -- (4.859,3.12);

%% ---------------- le kart, perpendiculaire à la piste -------------------------------------
\begin{scope}[shift={(I)}, rotate=\dev]
  \draw[line width=0.8pt, draw=solideB, fill=solideB!15, rounded corners=2pt]
    (-0.20,0.30) rectangle (1.72,0.50);
  \draw[line width=0.8pt, draw=solideB, fill=solideB!25, line join=round]
    (0.58,0.50) -- (1.02,0.50) -- (0.96,1.00) -- (0.66,1.00) -- cycle;   % dossier du siège
  \foreach \xr in {0,1.45} {
    \draw[line width=0.9pt, draw=black!80, fill=white] (\xr,0.24) circle (0.24);
    \fill[black!80] (\xr,0.24) circle (0.045); }
\end{scope}
\fill (I) circle (1.7pt);
\node[etiq, anchor=north west, inner sep=1.5pt] at (4.95,1.84) {$I$};

%% ---------------- la réponse (2) : l'angle phi et la conclusion --------------------------
\rappel{%
  \draw[line width=0.6pt, draw=solideA] ([shift=({90+\dev}:1.52)]I)
     arc ({90+\dev}:{90+\dev+\phif}:1.52);
  \node[etiq, text=solideA] at ([shift=({90+\dev+0.5*\phif}:1.70)]I) {$\varphi$};
  \draw[line width=0.6pt, draw=solideA] ([shift=({90+\dev-\phif}:1.52)]I)
     arc ({90+\dev-\phif}:{90+\dev}:1.52);
  \node[etiq, text=solideA] at ([shift=({90+\dev-0.5*\phif}:1.72)]I) {$\varphi$};
  \node[etiq, anchor=north west, text=solideA, align=left] at (0.25,4.20)
    {$\varphi = 36{,}9^\circ$\\la résultante du contact\\reste dans le cône :\\pas de glissement};}
\end{tikzpicture}
\end{document}
